\rhead{PH101: Physics}
\phantomsection 
\section*{Physics (PH101)}
\label{course:PH101}

\noindent \small{\hyperref[main:scheme]{$\leftarrow$ Back to Scheme}} \vspace{0.5cm}

\noindent \textbf{Credits:} 3 (3-0-0)

\subsection*{Course Content}
\noindent\textbf{Unit 1} \\
\textit{Quantum Mechanics (Computing Foundations):} Limitations of Classical Mechanics. Wave-Particle Duality, de Broglie’s Hypothesis, and Heisenberg Uncertainty Principle. \textbf{The Schrödinger Equation:} Time-Dependent and Time-Independent forms. Physical Interpretation of the Wave Function ($\Psi$). Applications: Particle in a 1D Infinite Potential Well (Energy Quantization) and \textbf{Quantum Tunneling} (Concept behind Flash Memory and tunneling limits in transistors). Introduction to Quantum Superposition and Qubits.

\vspace{0.3cm}

\noindent\textbf{Unit 2} \\
\textit{Semiconductor Physics (Hardware Foundations):} Band Theory of Solids: Conductors, Insulators, and Semiconductors. Intrinsic and Extrinsic Semiconductors (n-Type and p-Type), Fermi Levels. \textbf{Semiconductor Devices:} Physics of the P-N Junction, I-V Characteristics (Diode Logic), Zener Diodes, and Bipolar Junction Transistors (BJTs). Optoelectronics: Photodiodes, Solar Cells, and LEDs (Light Emitting Diodes).

\vspace{0.3cm}

\noindent\textbf{Unit 3} \\
\textit{Electromagnetism and Communication:} Electrostatics and Magnetostatics: Divergence and Curl of fields. \textbf{Maxwell’s Equations:} Integral and Differential forms and their physical significance (The logic of electromagnetic fields). EM Wave Propagation in Free Space. \textbf{Optical Fibers:} Total Internal Reflection, Numerical Aperture, and Signal Attenuation (Physics of the Internet backbone).

\vspace{0.3cm}

\noindent\textbf{Unit 4} \\
\textit{Oscillations and Signal Theory:} Simple Harmonic Motion (SHM). Damped Oscillations (Signal decay). Forced Oscillations and \textbf{Resonance} (Critical for clock speeds and filter design). Wave Motion: Transverse and Longitudinal waves, Principle of Superposition (Interference). Diffraction and Polarization (Basics of LCD screens).

\vspace{0.3cm}

\noindent\textbf{Unit 5} \\
\textit{Mechanics for Simulation (Game Physics):} Newton’s Laws and Galilean Invariance. \textbf{Kinematics and Dynamics:} Rectilinear and Projectile Motion (Trajectory algorithms). Work-Energy Theorem and Conservation of Energy. Impulse and Momentum: Conservation of Momentum during Collisions (Elastic/Inelastic) and Impact (Physics Engine logic). Introduction to Special Relativity: Time Dilation and Length Contraction (GPS Synchronization logic).

\newpage
\rhead{Curriculum Schema}